\section{Perspectiva geométrica: entender RoPE desde la frecuencia}

\subsection{Propiedad multiescala}

El decaimiento exponencial de $\theta_i$ implica que RoPE tiene una propiedad \textbf{multiescala}:

\begin{importantbox}{El espectro de RoPE}
\begin{itemize}
  \item \textbf{Índice bajo ($i$ pequeño)}: $\theta_i$ es grande, la frecuencia de rotación es alta $\to$ codifica dependencias de \textbf{corto alcance} (sintaxis, combinaciones de palabras como ``New York'')
  \item \textbf{Índice alto ($i$ grande)}: $\theta_i$ es pequeño, la frecuencia de rotación es baja $\to$ codifica dependencias de \textbf{largo alcance} (semántica, estructura de párrafo, referencias remotas)
\end{itemize}
\end{importantbox}

\subsection{La configuración de RoPE en Qwen3}

El RoPE estándar usa base = 10000, mientras que la serie Qwen3 usa base = 1000000. Una base mayor implica que todas las frecuencias son más bajas, la rotación del subespacio de baja frecuencia es más lenta, lo cual favorece capturar dependencias de mayor alcance y sustenta una context window más grande.

\subsection{Comparación con el esquema interleaved estándar}

Al presentar antes LLaMA se usó la versión interleaved ingenua de RoPE (que trata las dimensiones adyacentes como un subespacio 2D). Es posible que la implementación de Qwen3 difiera en la forma de dividir el subespacio; hace falta revisar el código concreto para confirmarlo.

\subsection{Resumen de la sección}

La esencia geométrica de RoPE consiste en aplicar rotaciones a distintas velocidades sobre los diferentes subespacios 2D del embedding. El subespacio de alta frecuencia se encarga de las relaciones sintácticas de corto alcance, y el subespacio de baja frecuencia se encarga de las relaciones semánticas de largo alcance. Qwen3 sustenta una ventana de contexto más larga aumentando la base.
